% 10-resolve.tex — Chapter 10: Name Resolution
\chapter{Name Resolution}
\label{chap:10-resolve}
\index{name resolution}
\index{sema\_resolve\_expr}

\pnum
Name resolution is implemented in \srcref{src/sema/resolve.h}.  It is Phase~1
of semantic analysis (\cref{sec:sema-twophase}): it visits every expression and
statement in a function body, populates \cname{->decl} pointers on identifier
nodes, and performs a set of source-to-source transforms before type inference
and checking.

% -------------------------------------------------------------------------
\section{Identifier resolution}
\label{sec:res-ident}
\index{identifier resolution}
% -------------------------------------------------------------------------

\pnum
\cname{sema\_resolve\_expr} handles \cname{EXPR\_IDENTIFIER} by:
\begin{enumerate}
\item Calling \cname{sema\_lookup(name)} to search locals then globals.
\item If found: setting \cname{e->decl = sym->decl}, \cname{e->type = sym->type},
  and \cname{e->is\_global = sym->is\_global}.
\item If not found: emitting an ``undefined identifier'' diagnostic and exiting.
\end{enumerate}

% -------------------------------------------------------------------------
\section{UFCS desugaring}
\label{sec:res-ufcs}
\index{UFCS}
\index{Uniform Function Call Syntax}
% -------------------------------------------------------------------------

\pnum
Uniform Function Call Syntax (UFCS) allows method-style calls to be written
as \lain{receiver.func(args)}.  The resolver detects \cname{EXPR\_CALL} whose
callee is \cname{EXPR\_MEMBER} and the member name resolves to a top-level
function.  It desugars the call by prepending the receiver as the first argument:

\begin{algorithm}
\textbf{UFCS desugaring} for \lain{recv.f(a, b)}:
\begin{enumerate}
\item Parse: \texttt{CALL(MEMBER(recv, "f"), [a, b])}
\item Lookup \texttt{"f"} in globals; if it resolves to a \cname{DECL\_FUNCTION} or
  \cname{DECL\_PROCEDURE}:
\item Rewrite to: \texttt{CALL(IDENT("f"), [recv, a, b])}
\end{enumerate}
\end{algorithm}

\pnum
If the member name does not resolve to a global function, the expression is
treated as a struct field access and left as \cname{EXPR\_MEMBER}.

% -------------------------------------------------------------------------
\section{Destructuring parameter expansion}
\label{sec:res-destruct}
\index{destructuring parameters}
% -------------------------------------------------------------------------

\pnum
When a function parameter carries a \cname{DECL\_DESTRUCT} node (syntax:
\lain{func f(\{x, y\} Point)}), the resolver:
\begin{enumerate}
\item Inserts a hidden parameter \cname{\_param\_N} with the struct type.
\item Resolves the struct type by scanning \cname{sema\_decls}.
\item For each destructured name, locates the corresponding field type and inserts
  a local binding mapping the name to \cname{\_param\_N.field}.
\end{enumerate}

\pnum
At emit time, the generated C function receives the hidden parameter, and each
destructured name is emitted as a field access on that hidden parameter.

% -------------------------------------------------------------------------
\section{Target platform constants}
\label{sec:res-platform}
\index{platform constants}
% -------------------------------------------------------------------------

\pnum
\srcref{src/sema/resolve.h} defines compile-time platform constants used to
evaluate \cname{@os} and \cname{@arch} builtin expressions:

\begin{center}
\begin{tabular}{lll}
\toprule
\textbf{Constant} & \textbf{Value} & \textbf{Condition} \\
\midrule
\cname{LAIN\_TARGET\_OS}   & \cname{LAIN\_OS\_LINUX}     & \texttt{\_\_linux\_\_} defined \\
\cname{LAIN\_TARGET\_OS}   & \cname{LAIN\_OS\_WINDOWS}   & \texttt{\_WIN32} defined \\
\cname{LAIN\_TARGET\_OS}   & \cname{LAIN\_OS\_MACOS}     & \texttt{\_\_APPLE\_\_} defined \\
\cname{LAIN\_TARGET\_ARCH} & \cname{LAIN\_ARCH\_X86\_64} & \texttt{\_\_x86\_64\_\_} defined \\
\cname{LAIN\_TARGET\_ARCH} & \cname{LAIN\_ARCH\_AARCH64} & \texttt{\_\_aarch64\_\_} defined \\
\bottomrule
\end{tabular}
\end{center}

\pnum
These constants are compiled into the Lain compiler binary and reflect the
host platform.  Cross-compilation would require passing them as flags; this
is not currently supported.

% -------------------------------------------------------------------------
\section{Scope building}
\label{sec:res-scopebuild}
\index{scope building}
% -------------------------------------------------------------------------

\pnum
\cname{sema\_build\_scope(decls, module\_path)} populates \cname{sema\_globals}
with all top-level declarations before Phase~1 begins.  It iterates over the
flat \cname{DeclList} and calls \cname{sema\_insert\_global} for each
\cname{DECL\_FUNCTION}, \cname{DECL\_PROCEDURE}, \cname{DECL\_STRUCT},
\cname{DECL\_ENUM}, and \cname{DECL\_TYPE\_ALIAS}.

\pnum
Global scope is built once per module before any function body is resolved.
This allows forward references: a function \texttt{f} can call a function
\texttt{g} that appears later in the source file.
